\section{版式样张说明}\label{sec:sample-intro}

本节文字只用于检查双栏正文的字号、行距、中文换行和段落缩进。正式稿将依据已经确认的内容蓝图分节写作，不会保留这些占位段落。样张引用两篇已核验文献，以验证 GB/T 7714—2015 顺序编码、超链接和参考文献悬挂缩进是否正常：FedAvg 奠定了本地更新与周期聚合的基本范式\cite{McMahan2017FedAvg}，移动边缘网络综述则提供了无线资源和设备异构的研究背景\cite{Lim2020MobileEdge}。

双栏版式的主要目的，是在有限页面内建立稳定的阅读节奏，而不是简单压缩内容。正文应保持两字符首行缩进，中文标点不能出现在不合适的行首位置，英文术语和数学符号也不应挤压相邻文字。本段故意混排 federated learning、AirComp、\(B_{k,t}\) 等字符，用于观察中英文基线、字距和断行效果。

\subsection{标题层级与交叉引用}\label{subsec:sample-crossref}

一级标题、二级标题和正文之间需要形成清晰但克制的层级。后文将通过\cref{eq:sample-risk}、\cref{tab:sample-layout}和\cref{fig:sample-placeholder}检查公式、表格和图片的自动编号；这些编号均由 LaTeX 生成，不在正文中手工录入。引用中文综述\cite{Zhang2023EdgeFLSurvey}还可同时检查中文作者、期刊名和 DOI 的排版。

\section{原生公式与变量说明}\label{sec:sample-equation}

公式应服务于概念界定，并在前后给出解释。设系统包含 \(K\) 个终端，第 \(k\) 个终端拥有数据集 \(\mathcal D_k\)，样本数为 \(n_k\)，模型参数为 \(\bm w\)，则全局经验风险的样张公式写为
\begin{equation}\label{eq:sample-risk}
  \min_{\bm w} F(\bm w)
  = \sum_{k=1}^{K} \frac{n_k}{N} F_k(\bm w),
  \qquad N=\sum_{k=1}^{K}n_k .
\end{equation}
式中，\(F_k(\bm w)\) 表示终端 \(k\) 的局部经验风险，\(N\) 为总样本数。该公式采用原生数学环境排版，可被自动编号、检索和交叉引用，不依赖 PNG、SVG 或公式截图。

\subsection{三线表样张}\label{subsec:sample-table}

\cref{tab:sample-layout}用于验证单栏表格的宽度、三线表样式、表题位置和表下注释。正式稿中的比较表将根据内容宽度决定使用单栏表或跨双栏表，不通过缩小到难以阅读的字号来强行塞入。

\begin{table}[t]
  \centering
  \caption{版式组件检查表}
  \label{tab:sample-layout}
  \begin{threeparttable}
    \small
    \begin{tabularx}{\columnwidth}{@{}lYYY@{}}
      \toprule
      组件 & 编号 & 引用 & 样张状态 \\
      \midrule
      原生公式 & 自动 & \texttt{eq:} & 已配置 \\
      三线表 & 自动 & \texttt{tab:} & 已配置 \\
      TikZ 图 & 自动 & \texttt{fig:} & 临时占位 \\
      参考文献 & 顺序 & \texttt{cite} & GB/T 7714 \\
      \bottomrule
    \end{tabularx}
    \begin{tablenotes}[flushleft]\footnotesize
      \item 注：该表仅测试排版功能，不包含研究数据。
    \end{tablenotes}
  \end{threeparttable}
\end{table}


表格表明，各类对象应采用与其功能相匹配的排版环境：公式需要可引用的数学编号，图片需要明确题注，参考文献则必须由数据库生成顺序编号。这里的“状态”仅指版式测试结果，不代表论文内容已经完成。

\subsection{临时占位图样张}\label{subsec:sample-figure}

\cref{fig:sample-placeholder}是一幅项目内原创的 TikZ 临时占位图，仅用于检查黑白矢量线条、图宽、图题和浮动位置。它不会进入最终正文，后续将由经过内容审查的系统模型图替换。

\begin{figure}[t]
  \centering
  \resizebox{0.96\columnwidth}{!}{%
    \begin{tikzpicture}[
      node distance=10mm,
      box/.style={draw=black, rounded corners=1.5pt, minimum width=19mm,
                  minimum height=8mm, align=center, font=\small, fill=black!4},
      arr/.style={-{Stealth[length=4pt,width=5pt]}, line width=0.6pt}
    ]
      \node[box] (device) {无线终端};
      \node[box, right=of device] (edge) {边缘节点};
      \node[box, right=of edge] (cloud) {云服务器};
      \draw[arr] ([yshift=1.2mm]device.east) -- node[above,font=\scriptsize] {局部更新} ([yshift=1.2mm]edge.west);
      \draw[arr] ([yshift=-1.2mm]edge.west) -- node[below,font=\scriptsize] {模型下发} ([yshift=-1.2mm]device.east);
      \draw[arr] ([yshift=1.2mm]edge.east) -- node[above,font=\scriptsize] {边缘聚合} ([yshift=1.2mm]cloud.west);
      \draw[arr] ([yshift=-1.2mm]cloud.west) -- node[below,font=\scriptsize] {全局模型} ([yshift=-1.2mm]edge.east);
    \end{tikzpicture}%
  }
  \caption{版式占位图：端—边—云消息路径（作者绘制，仅用于样张）}
  \label{fig:sample-placeholder}
\end{figure}


图后说明用于检验“图前提出问题—图中表达关系—图后解释作用”的写法。当前图只展示设备、边缘与云之间的消息方向，不表达性能比较，也不含来自外部论文的图形元素。

\section{双栏分页与参考文献样张}\label{sec:sample-pagination}

本节增加适量占位文字，以检查页面是否能稳定生成二至三页，并观察浮动体与参考文献在双栏中的分页行为。正式论文不会为了填满页面而保留无信息段落，而会依据每节字数预算安排论证密度。样张文本应足以暴露明显的栏宽、页边距、孤行和标题悬空问题。

排版检查还应区分可自动验证和必须人工判断的项目。字体是否嵌入、页面尺寸、未定义引用和 overfull box 可以通过命令检查；标题是否显得拥挤、摘要是否过长、图表与正文是否形成解释关系，则需要查看渲染结果。两类检查缺一不可。

页眉仅保留论文简称和页码，正文没有装饰性边框、彩色背景或水印。黑白版式可减少打印差异，并为后续机制图和比较表提供一致的视觉基础。图片、表格和公式均设置稳定标签，后续移动位置时不需要修改正文中的编号。

参考文献样张只放入三条已核验文献，用于测试 Biber 和 GB/T 7714—2015 样式；它不代表最终文献库缩减。完整论文仍应使用已确认的 43 篇文献，并在正文完成后执行引文与文献表的双向覆盖审计。

为避免紧窄栏内出现明显溢出，长英文题名和 DOI 应允许合理换行，表格文字应自动折行，公式过长时应改用 aligned、split 或 multline 环境。最终工程还需在每次新增复杂图表后重新执行相同检查，而不能只依赖本次样张通过。

本阶段的交付边界是版式验证。样张确认后，下一阶段才能按照内容蓝图迁移经核验的题录、章节和引用；在获得确认之前，不会把 v1.2 正文批量复制到这些章节文件中。
